\subsection{两模态梁的领先阶非退化性}

第\ref{sec:experiments}节表明，弱非线性梁的长期预测误差主要受制于相位推进。要从几何结构上判断小微扰或离散下拟周期运动能否持续，需引入经典KAM定理\citep[第1.e节，The Classical KAM Theorem]{Poschel2001}（定理及证明归属外部文献），其核心前提是可积参考系统满足Kolmogorov频率非退化条件。然而，对于式\eqref{eq:learning-model}的两模态梁，谐振子参考作用量$I_j=(X_j^2+Y_j^2)/2$下的二次主部$h_0=\nu^TI$频率恒为$\nu$、满足$\nabla_I^2h_0=0$，本身完全退化；非退化性只能来自静态余弦磁势的高阶项。为此，本文对余弦磁势作四次展开并沿$X_j=\sqrt{2I_j}\cos\theta_j$角平均，导出领先阶正规形项：
\begin{equation}
\begin{aligned}
\Delta h_{\mathrm{phys}}&=-\frac{(d^TX)^4}{24}+O((d^TX)^6),\\
\overline h_4(I)&=-\frac{d_1^4I_1^2+4d_1^2d_2^2I_1I_2+d_2^4I_2^2}{16}.
\end{aligned}
\label{eq:beam-quartic-average}
\end{equation}
因四次项非零活跃谐波$0<|k|_1\leq4$的最小分母$|k\cdot\nu|=2$远离共振，式\eqref{eq:beam-quartic-average}给出非共振正规形领先项，其作用量Hessian及行列式为
\begin{equation}
B_4=-\begin{pmatrix}
d_1^4/8&d_1^2d_2^2/4\\
d_1^2d_2^2/4&d_2^4/8
\end{pmatrix},\qquad
\det B_4=-\frac{3d_1^4d_2^4}{64}.
\label{eq:beam-leading-twist}
\end{equation}
代入本文$d$得$\det B_4=-1.30371\times10^{-6}\neq0$：磁势确实提供了领先阶频率非退化性，但这尚不构成完整系统在作用量域上的非退化下界；高阶正规形余项仍待控制，且谐振子参考量$I_j$在全系统中并非精确守恒量。

\subsection{连续频率与中点步长}

为在学习幅值区内检验上述预测，取表\ref{table:beam_kam}三组半径及固定初相位$(0.37,1.13)$，置$X_j(0)=r_j\cos\theta_j$、$Y_j(0)=-r_j\sin\theta_j$，以DOP853（容差$10^{-12}$、$10^{-14}$）积分至$\tau=1023.9809$（$89.6869\,\mathrm s$，采样间隔$\Delta\tau=0.0399993$）。图\ref{fig:beam_kam}左图以FFT（分辨率约$0.00614$）辨认双主峰，微小频移则由$X_j-iY_j$相位增量的光滑加权平均提取：中心工况加密容差频率变化小于$7\times10^{-13}$，分窗与降采样检验确认表中频移均高于$10^{-10}$分辨阈值。在中心初值处对初值参数$I_j^0=r_j^2/2$作两级中心差分，得频率响应矩阵
\[
\frac{\partial\omega}{\partial I^0}\approx10^{-3}
\begin{pmatrix}
-0.32411&-1.31818\\
-1.31818&-1.34011
\end{pmatrix},
\]
其行列式$-1.30326\times10^{-6}$及各分量与$B_4$的最大差不超过$3.32\times10^{-7}$，与四次平均预测吻合。该矩阵是在固定初相位下对参考初值参数的响应，并非精确作用量坐标中的频率梯度；上述经验分辨阈值也不是严格的频率误差界。

进而将训练基准步长$h_0=0.00999981$依次减半至$h_0/32$积分至同一终点（图\ref{fig:beam_kam}右图）：两模态频率误差观察阶为$1.9994$--$2.0000$，但中心工况在$h_0$下的第二模态离散频差（约$1.94\times10^{-3}$）高达其物理频移（$1.62\times10^{-5}$）的120倍，在所测工况中，步长缩减逾一个数量级后，第二模态离散频差才降至物理频移之下。这与第\ref{sec:experiments}节中积分误差主导连续误差的现象一致。

\subsection{KAM条件与学习误差}

将频率诊断联系到不变环面的持续性，还需满足相应定理的完整前提。\citet[定理1]{Shang1999}针对解析、可积且Kolmogorov非退化的参考系统，说明相容解析辛格式在足够小步长下保持大部分非共振不变环面。该定理及证明归属原作者；本文完整两模态能量的可积性尚未建立，领先阶正规形也仍需连同高阶余项检验，因而不能仅由$\det B_4\neq0$或频率收敛直接援用其结论。

\citet[定理11，假设2--5]{Chen2022}在目标系统可积、非退化等条件下，将规定的泛化界与KAM阈值联系起来。该文采用$p>2M$的梯度$L^p$损失，两自由度时为$p>4$；$p$是积分指数，不是导数阶数。其统一能量误差估计先利用Hamiltonian常数平移不改变动力学，将能量误差归一化为零均值，再结合Poincaré与Sobolev不等式，以$W^{1,p}$控制能量的统一界。目标Hamiltonian及网络激活函数的$C^\infty$正则性是另一项独立假设。该文假设1还要求数据分布密度在所考察区域内有正下界，并满足所述泛化误差界；这些条件及阈值常数尚未在本文训练协议中核验。有限轨迹上的均方损失因此不能直接替代该统一界。

本节已计算领先阶$\det B_4\neq0$、双主频有限阶筛查（$0<|k|_1\leq32$内最小$|k\cdot\omega|\approx0.15215$）及离散频移收敛阶。全系统可积性、高阶正规形余项、作用量域上的非退化下界、全整数格Diophantine条件及学习误差所需的统一界仍未认证，现有证据不构成完整系统或已训练网络的严格不变环面证明。
